\section*{Executive summary}
\addcontentsline{toc}{section}{Executive summary}

The \network{} operated eight instrumented towers along 140 kilometres of
low lying dune and saltmarsh coast throughout 2024, recording turbulent fluxes
of heat, moisture, and carbon dioxide at 20 hertz alongside a standard
meteorological suite at one minute resolution. This report documents the
network as built, the processing chain applied to the raw records, the quality
control outcomes, and the annual flux budget that follows.

Four findings are reported.

Data return across the network was 91.4 percent after quality control, against
a design target of 90 percent. The shortfall at the two southernmost towers,
MC06 and MC07, is attributable to salt contamination of the open path gas
analyser windows and is the subject of Recommendation~\ref{rec:enclosure}. The
same two towers account for 61 percent of all records rejected for a failed
stationarity test.

The saltmarsh sites were a net carbon sink of $-214 \pm 31$ grams of carbon per
square metre over the year, while the dune sites were approximately neutral at
$-12 \pm 26$. The difference is driven almost entirely by the growing season,
between day 120 and day 260, and the two land cover classes are
indistinguishable outside it.

Nocturnal fluxes remain the dominant source of uncertainty. Under the friction
velocity filter adopted in Section~\ref{sec:qc-ustar}, 38 percent of nocturnal
half hours are rejected and gap filled, and the choice of threshold between
0.08 and 0.16 metres per second moves the annual budget by 44 grams of carbon
per square metre. This is larger than the reported uncertainty on either land
cover class and is stated separately rather than folded into it.

The sea breeze circulation produces a systematic advective term that the eddy
covariance method does not capture. Section~\ref{sec:results-advection}
estimates it at 8 to 14 percent of the midday carbon flux on days with an
onshore wind reversal, which occurred on 96 days of the year. The network as
currently instrumented cannot measure this term, and
Recommendation~\ref{rec:profile} proposes the vertical profile system that
would.
